\chapter{Итоговые результаты и сложность оптимизированной программы}

\section{Сводные результаты}

Сводные результаты всех шагов приведены в таблице~\ref{tab:final-cycles}
и на рисунках~\ref{fig:steps-compress} и~\ref{fig:steps-decompress}. Это
медианы пяти раундов сценария \texttt{bench/steps.sh}, который собрал
программу каждого шага из её коммита и измерил все шаги подряд на одном
ядре процессора. Столбец <<к пред.>> "--- ускорение относительно предыдущего
шага, <<к исх.>> "--- относительно исходной программы; для сборки с PGO
предыдущим шагом считается сборка только с LTO.

\begin{table}[H]
\caption{Такты процессора на весь корпус по шагам, млрд}
\label{tab:final-cycles}
\begin{tabletext}
\centering
\setlength{\tabcolsep}{4pt}
\begin{tabular}{|l|r|r|r|r|r|r|r|r|}
\hline
\multirow{3}{*}{\textbf{Шаг}} & \multicolumn{4}{c|}{\textbf{debian (Ryzen 7 5700U)}} & \multicolumn{4}{c|}{\textbf{hppavillion (i3-1115G4)}} \\
\cline{2-9}
 & \multicolumn{2}{c|}{\textbf{сжатие}} & \multicolumn{2}{c|}{\textbf{распаковка}} & \multicolumn{2}{c|}{\textbf{сжатие}} & \multicolumn{2}{c|}{\textbf{распаковка}} \\
\cline{2-9}
 & \textbf{такты} & \textbf{к пред.} & \textbf{такты} & \textbf{к пред.} & \textbf{такты} & \textbf{к пред.} & \textbf{такты} & \textbf{к пред.} \\
\hline
00 исходная & 41,85 & & 7,57 & & 42,15 & & 7,38 & \\
01 блочный в/в & 39,93 & 1,05 & 4,64 & \textbf{1,63} & 41,68 & 1,01 & 4,42 & \textbf{1,67} \\
02 встраивание & 41,94 & 0,95 & 3,24 & \textbf{1,43} & 40,86 & 1,02 & 3,42 & \textbf{1,29} \\
03 без сдвига & 27,62 & \textbf{1,52} & 3,25 & 1,00 & 31,65 & \textbf{1,29} & 3,36 & 1,02 \\
04 по 8 байт & 16,65 & \textbf{1,66} & 3,24 & 1,00 & 17,99 & \textbf{1,76} & 3,38 & 0,99 \\
05 AVX-512BW & 16,29 & 1,02 & 3,26 & 0,99 & 14,90 & \textbf{1,21} & 3,37 & 1,00 \\
06 AVX2 & 14,40 & \textbf{1,13} & 3,34 & 0,98 & 14,69 & 1,01 & 3,40 & 0,99 \\
07 LTO & 14,19 & 1,01 & 3,32 & 1,00 & 14,66 & 1,00 & 3,33 & 1,02 \\
07 LTO + PGO & 13,47 & \textbf{1,05} & 3,14 & \textbf{1,06} & 14,80 & 0,99 & 3,10 & \textbf{1,07} \\
\hline
\textbf{Итого к исходной} & & \textbf{3,11} & & \textbf{2,41} & & \textbf{2,85} & & \textbf{2,38} \\
\hline
\end{tabular}
\end{tabletext}
\end{table}

\begin{figure}[H]
\centering
\begin{tikzpicture}
\begin{axis}[
  ybar, width=\textwidth, height=7.5cm, bar width=7pt,
  ylabel={Сжатие корпуса, млрд тактов}, ymin=0, ymax=48,
  symbolic x coords={00,01,02,03,04,05,06,07a,07b}, xtick=data,
  xticklabels={00,01,02,03,04,05,06,{07 LTO},{07 PGO}},
  xticklabel style={font=\small},
  xlabel={Шаг оптимизации},
  nodes near coords, nodes near coords style={font=\tiny, rotate=90, anchor=west},
  legend style={at={(0.98,0.98)}, anchor=north east, font=\small},
  enlarge x limits=0.07, grid=major, grid style={gray!25},
]
\addplot[fill=gray!40] table[x=step, y=compress] {data/steps-debian.dat};
\addlegendentry{debian}
\addplot[fill=gray!85] table[x=step, y=compress] {data/steps-hppavillion.dat};
\addlegendentry{hppavillion}
\end{axis}
\end{tikzpicture}
\caption{Такты сжатия всего корпуса по шагам оптимизации}
\label{fig:steps-compress}
\end{figure}

\begin{figure}[H]
\centering
\begin{tikzpicture}
\begin{axis}[
  ybar, width=\textwidth, height=7cm, bar width=7pt,
  ylabel={Распаковка корпуса, млрд тактов}, ymin=0, ymax=9,
  symbolic x coords={00,01,02,03,04,05,06,07a,07b}, xtick=data,
  xticklabels={00,01,02,03,04,05,06,{07 LTO},{07 PGO}},
  xticklabel style={font=\small},
  xlabel={Шаг оптимизации},
  nodes near coords, nodes near coords style={font=\tiny, rotate=90, anchor=west},
  legend style={at={(0.98,0.98)}, anchor=north east, font=\small},
  enlarge x limits=0.07, grid=major, grid style={gray!25},
]
\addplot[fill=gray!40] table[x=step, y=decompress] {data/steps-debian.dat};
\addlegendentry{debian}
\addplot[fill=gray!85] table[x=step, y=decompress] {data/steps-hppavillion.dat};
\addlegendentry{hppavillion}
\end{axis}
\end{tikzpicture}
\caption{Такты распаковки всего корпуса по шагам оптимизации}
\label{fig:steps-decompress}
\end{figure}

\screenshot{steps-debian}{Результат сценария bench/steps.sh на компьютере debian}{fig:steps-debian}{%
\texttt{cat bench/log/debian/steps/steps.txt} в терминале шириной не менее
100 символов. Повторить измерение: \texttt{ROUNDS=5 CPU=15 bench/steps.sh}.}

\screenshot{steps-hpp}{Результат сценария bench/steps.sh на компьютере hppavillion}{fig:steps-hpp}{%
\texttt{cat bench/log/hppavillion/steps/steps.txt}. Повторить измерение:
\texttt{ROUNDS=5 CPU=3 bench/steps.sh} (ноутбук подключён к сети и не нагружен).}

В~секундах на компьютере debian сжатие всего корпуса ускорилось с 10,10 до
3,29~с (в 3,07~раза, с 28,9 до 88,7~МБ/с), распаковка "--- с 2,07 до
0,90~с (в 2,29~раза, со 141 до 324~МБ/с), таблица~\ref{tab:wall} и
рисунок~\ref{fig:throughput}.

\begin{table}[H]
\caption{Время обработки всего корпуса на компьютере debian (сумма медиан по файлам, 9 запусков)}
\label{tab:wall}
\begin{tabletext}
\centering
\begin{tabular}{|l|r|r|r|r|}
\hline
\textbf{Шаг} & \textbf{Сжатие, с} & \textbf{Распаковка, с} & \textbf{Сжатие, МБ/с} & \textbf{Распаковка, МБ/с} \\
\hline
00 исходная & 10,10 & 2,07 & 28,9 & 141,2 \\
01 блочный ввод-вывод & 10,27 & 1,24 & 28,4 & 235,3 \\
02 встраивание & 9,89 & 0,95 & 29,5 & 306,8 \\
03 без сдвига буфера & 7,04 & 0,93 & 41,5 & 313,1 \\
04 сравнение по 8 байт & 3,88 & 0,93 & 75,3 & 314,6 \\
05 AVX-512BW & 3,92 & 0,92 & 74,5 & 316,4 \\
06 AVX2 & 3,48 & 0,94 & 83,8 & 309,4 \\
07 LTO + PGO & 3,29 & 0,90 & 88,7 & 323,9 \\
\hline
\end{tabular}
\end{tabletext}
\end{table}

\begin{figure}[H]
\centering
\begin{tikzpicture}
\begin{axis}[
  width=0.9\textwidth, height=7cm,
  symbolic x coords={00,01,02,03,04,05,06,07b}, xtick=data,
  xticklabels={00,01,02,03,04,05,06,07},
  xlabel={Шаг оптимизации}, ylabel={Скорость, МБ/с},
  ymin=0, ymax=350, grid=major, grid style={gray!25},
  legend pos=north west, legend style={font=\small},
  nodes near coords, nodes near coords style={font=\tiny, anchor=south},
]
\addplot[mark=*, thick] table[x=step, y=dmbs] {data/wall-debian.dat};
\addlegendentry{распаковка}
\addplot[mark=square*, thick, dashed] table[x=step, y=cmbs] {data/wall-debian.dat};
\addlegendentry{сжатие}
\end{axis}
\end{tikzpicture}
\caption{Скорость сжатия и распаковки по шагам, debian}
\label{fig:throughput}
\end{figure}

Ускорение по отдельным файлам (рисунок~\ref{fig:per-file-total})
составляет от 2 до 4,4~раза при сжатии и от 1,2 до 3,8~раза при
распаковке. Сильнее всего ускорились random.dat (больше всего токенов на
байт) и mixed.dat (самые длинные цепочки совпадений), меньше всего "---
zeros.dat, который и в исходной программе обрабатывался быстрее всех.

\begin{figure}[H]
\centering
\begin{tikzpicture}
\begin{axis}[
  ybar, width=\textwidth, height=7.5cm, bar width=5pt,
  ylabel={Ускорение, раз}, ymin=0, ymax=4.8,
  xtick=data, xticklabels from table={data/perfile-debian.dat}{file},
  xticklabel style={rotate=60, anchor=east, font=\small\ttfamily},
  legend style={at={(0.98,0.98)}, anchor=north east, font=\small},
  enlarge x limits=0.04, grid=major, grid style={gray!25},
]
\addplot[fill=gray!40] table[x=index, y=stotal] {data/perfile-debian.dat};
\addlegendentry{сжатие}
\addplot[fill=gray!85] table[x=index, y=dtotal] {data/perfile-debian.dat};
\addlegendentry{распаковка}
\end{axis}
\end{tikzpicture}
\caption{Итоговое ускорение по файлам, шаг 07 против исходной программы (debian, время)}
\label{fig:per-file-total}
\end{figure}

На всех шагах и на обоих компьютерах архивы всех 16 файлов корпуса
совпадают байт в байт с архивами исходной программы
(рисунок~\ref{fig:sha256}), коэффициент сжатия корпуса остаётся равным
0,4753, а набор тестов \texttt{make check} проходит
(рисунок~\ref{fig:make-check}).

\screenshot{sha256}{Совпадение контрольных сумм архивов исходной и оптимизированной программ}{fig:sha256}{%
\texttt{diff bench/log/debian/00-baseline/raw/archives.sha256
bench/log/debian/07-lto-pgo/raw/archives.sha256 \&\& echo identical},
а также \texttt{head -5} одного из файлов, чтобы были видны суммы.}

\screenshot{make-check}{Успешное выполнение набора тестов make check}{fig:make-check}{%
На текущей версии: \texttt{make check}; снять итоговую сводку
<<Testsuite summary>> с \texttt{PASS} и \texttt{FAIL: 0}.}

\section{Профиль оптимизированной программы}

Метрики процессора на файле dickens для всех шагов показаны на
рисунке~\ref{fig:profile-steps}, сравнение исходной и оптимизированной
программ "--- в таблице~\ref{tab:final-profile}. Профильная сборка на всех
шагах одинакова (\texttt{-O2 -g}, без LTO и PGO), поэтому шаг~07 в
профиле совпадает с шагом~06. На debian профильная сборка шагов 05--07
медленнее шага~04 на dickens: это один из файлов, на которых реализация
AVX2 проигрывает сравнению по 8 байт (раздел~\ref{sec:avx2}).

\begin{figure}[H]
\centering
\begin{tikzpicture}
\begin{axis}[
  width=0.9\textwidth, height=7cm,
  symbolic x coords={00,01,02,03,04,05,06,07b}, xtick=data,
  xticklabels={00,01,02,03,04,05,06,07},
  xlabel={Шаг оптимизации}, ylabel={Сжатие dickens, млн},
  ymin=0, ymax=3000, grid=major, grid style={gray!25},
  legend style={at={(0.98,0.98)}, anchor=north east, font=\small},
]
\addplot[mark=*, thick] table[x=step, y=cinstr] {data/profile-debian.dat};
\addlegendentry{команды, debian}
\addplot[mark=o, thick, dashed] table[x=step, y=cinstr] {data/profile-hppavillion.dat};
\addlegendentry{команды, hppavillion}
\addplot[mark=square*, thick] table[x=step, y=ccycles] {data/profile-debian.dat};
\addlegendentry{такты, debian}
\addplot[mark=square, thick, dashed] table[x=step, y=ccycles] {data/profile-hppavillion.dat};
\addlegendentry{такты, hppavillion}
\end{axis}
\end{tikzpicture}
\caption{Команды и такты при сжатии файла dickens по шагам (профильная сборка)}
\label{fig:profile-steps}
\end{figure}

\begin{table}[H]
\caption{Профиль исходной и оптимизированной программ, dickens, debian}
\label{tab:final-profile}
\begin{tabletext}
\centering
\begin{tabular}{|l|r|r|r|}
\hline
\textbf{Метрика} & \textbf{Шаг 00} & \textbf{Шаг 07} & \textbf{Изменение} \\
\hline
Сжатие: такты, млн & 1150 & 553 & в 2,08 раза меньше \\
Сжатие: команды, млн & 2747 & 1099 & в 2,50 раза меньше \\
Сжатие: ветвления, млн & 541 & 194 & в 2,79 раза меньше \\
Сжатие: ошибки предсказания, млн & 15,8 & 7,6 & в 2,09 раза меньше \\
Распаковка: такты, млн & 350 & 170 & в 2,06 раза меньше \\
Распаковка: команды, млн & 862 & 346 & в 2,49 раза меньше \\
Cachegrind, сжатие и распаковка: команды, млн & 3642 & 1441 & в 2,53 раза меньше \\
Cachegrind: промахи L1D при чтении, млн & 14,1 & 12,8 & на 10~\% меньше \\
Cachegrind: ошибки предсказания ветвлений, млн & 34,2 & 15,9 & в 2,16 раза меньше \\
\hline
\end{tabular}
\end{tabletext}
\end{table}

Плоский профиль gprof оптимизированной программы
(рисунок~\ref{fig:final-gprof}) показывает, что теперь время делится между
циклом сжатия \texttt{compress\_avx2} (32~\%) и сравнением кандидатов
\texttt{match\_length\_avx2} (26~\%); распаковка занимает 8~\%, вывод
байта \texttt{output\_byte} "--- 7~\%, хеш "--- 6~\%. Вызовы библиотеки
ввода-вывода и копирование памяти из профиля исчезли. Число промахов кеша
почти не изменилось: все структуры данных программы (менее 1~МиБ)
по-прежнему помещаются в кеш последнего уровня~\cite{drepper}, где
Cachegrind насчитывает лишь несколько тысяч промахов, и оставшееся время "--- это вычисления и
зависимые переходы по цепочкам хеш-таблицы.

\screenshot{final-gprof}{Плоский профиль gprof оптимизированной программы}{fig:final-gprof}{%
\texttt{head -20 bench/log/debian/07-lto-pgo/raw/profile/gprof.txt}; при
желании рядом \texttt{perf report} по строкам:
\texttt{head -20 bench/log/debian/07-lto-pgo/raw/profile/perf-lines.txt}.}

\section{Оценка алгоритмической сложности оптимизированной программы}

\textbf{Сжатие.} Для каждого токена оптимизированная программа выполняет:
\begin{enumerate}
\item вычисление хеша и выбор головы цепочки "--- $O(1)$;
\item просмотр не более $C$ кандидатов; сравнение каждого занимает не более
$\lceil M/w \rceil$ шагов, где $w$ "--- ширина сравнения: 8~байт в
переносимой версии, 32~байта с AVX2, 64~байта с AVX-512, "--- $O(C \cdot
\lceil M/w \rceil)$;
\item запись токена в блок архива "--- $O(1)$;
\item вставку в хеш-таблицу пройденных позиций "--- $O(1)$ на байт входа.
\end{enumerate}
Сдвига буфера на каждый токен больше нет; перенос данных при дополнении
общего буфера стоит $O(n)$ на весь файл (не более 0,15~байта на входной
байт). Итого
\begin{equation}
\label{eq:final}
T_{\text{сж}}(n) = O\bigl(n + T \cdot C \cdot \lceil M/w \rceil\bigr) =
O\bigl(n \cdot C \cdot \lceil M/w \rceil\bigr).
\end{equation}
Асимптотически сжатие, как и прежде, линейно: $T_{\text{сж}}(n) = O(n)$.
Изменилась константа: из стоимости токена исключено слагаемое $L = 4096$, а
член поиска уменьшился в $w$ раз; в худшем случае "--- с $C \cdot M + L
\approx 70\,000$ до $C \cdot \lceil M/32 \rceil \approx 2300$ шагов на байт
при AVX2. Улучшить асимптотику ниже $O(n)$ невозможно, поскольку каждый
входной байт нужно хотя бы прочитать.

\textbf{Распаковка.} $T_{\text{расп}}(n) = O(n)$, как и прежде; каждый байт
теперь обходится записью в массив и сдвигом битового буфера вместо вызовов
\texttt{fgetc} и \texttt{fputc}.

\textbf{Память.} Сжатие использует таблицу голов (512~КиБ), массив
\texttt{ht\_next} (32~КиБ), общий буфер данных (64~КиБ) и блок
побитового писателя (64~КиБ), всего около 672~КиБ против 552~КиБ в
исходной программе; распаковка вместо одного окна в 4~КиБ использует ещё
два блока по 64~КиБ. Объём памяти по-прежнему $O(1)$ относительно размера
файла и увеличился не более чем на 128~КиБ.

\textbf{Экспериментальная проверка.} Префиксы тех же синтетических файлов
размером от 1 до 20~МБ сжимались исходной и оптимизированной
программами (таблица~\ref{tab:ipb}). Обе
зависимости линейны: число команд на байт входа не зависит от размера
входа, а у оптимизированной программы оно меньше в 1,7--4,2~раза. Сильнее
всего константа уменьшилась на случайных данных, где в исходной программе
преобладал сдвиг буфера.

\begin{table}[H]
\caption{Команды процессора на байт входа при сжатии (hppavillion)}
\label{tab:ipb}
\begin{tabletext}
\centering
\begin{tabular}{|l|r|r|r|}
\hline
\textbf{Файл} & \textbf{Исходная} & \textbf{Оптимизированная} & \textbf{Отношение} \\
\hline
random.dat & 384 & 92 & 4,2 \\
text.dat & 251 & 87 & 2,9 \\
zeros.dat & 33 & 20 & 1,7 \\
\hline
\end{tabular}
\end{tabletext}
\end{table}
